\chapter{Results and Discussions}\label{ch:4}

This chapter presents the analysis and interpretation of the data gathered from the respondents regarding the developed web-based learning tool, \textit{StatiCalcs}. The findings focus on the respondents’ perceptions in terms of usability, accessibility, and satisfaction, as well as their preference for different learning resources in studying Statics of Rigid Bodies. Frequency distribution, weighted mean, and independent samples t-test were utilized to analyze the gathered data.

\section{Level of Perception of Students Regarding StatiCalcs in Terms of Usability}

\begin{table}[H]
\centering
\caption{Frequency Distribution of Students’ Responses on Usability}
\renewcommand{\arraystretch}{1.3}
\begin{tabular}{p{8cm}ccccc}
\hline
\textbf{Statement} & \textbf{SA} & \textbf{A} & \textbf{N} & \textbf{D} & \textbf{SD}\\
\hline
The StatiCalcs website is easy to navigate. & 155 & 48 & 0 & 0 & 0\\
I can complete calculations quickly using the website. & 140 & 54 & 0 & 2 & 0\\
The website provides clear feedback after performing calculations. & 150 & 52 & 0 & 1 & 0\\
I can easily correct mistakes when entering inputs in the calculators. & 159 & 42 & 0 & 0 & 0\\
Overall, StatiCalcs is user-friendly. & 160 & 43 & 0 & 0 & 0\\
\hline
\end{tabular}
\end{table}

\begin{table}[H]
\centering
\caption{Weighted Mean of Students’ Perception on Usability}
\renewcommand{\arraystretch}{1.3}
\begin{tabular}{p{8cm}cc}
\hline
\textbf{Statement} & \textbf{Weighted Mean} & \textbf{Interpretation}\\
\hline
The StatiCalcs website is easy to navigate. & 4.69 & Strongly Agree\\
I can complete calculations quickly using the website. & 4.57 & Strongly Agree\\
The website provides clear feedback after performing calculations. & 4.66 & Strongly Agree\\
I can easily correct mistakes when entering inputs in the calculators. & 4.70 & Strongly Agree\\
Overall, StatiCalcs is user-friendly. & 4.71 & Strongly Agree\\
\textbf{Overall Weighted Mean} & \textbf{4.66} & \textbf{Strongly Agree}\\
\hline
\end{tabular}
\end{table}

The findings reveal that the student respondents strongly agreed that \textit{StatiCalcs} is usable and user-friendly. The statement, ``Overall, StatiCalcs is user-friendly,'' obtained the highest weighted mean of 4.71, indicating a highly positive evaluation of the platform’s overall usability. Meanwhile, the statement, ``I can complete calculations quickly using the website,'' obtained the lowest weighted mean of 4.57. Despite this, all indicators remained within the ``Strongly Agree'' interpretation range. The results suggest that the respondents perceived the developed web-based learning tool as easy to navigate, efficient, and beneficial for solving Statics problems.

\section{Level of Perception of Students Regarding StatiCalcs in Terms of Accessibility}

\begin{table}[H]
\centering
\caption{Frequency Distribution of Students’ Responses on Accessibility}
\renewcommand{\arraystretch}{1.3}
\begin{tabular}{p{8cm}ccccc}
\hline
\textbf{Statement} & \textbf{SA} & \textbf{A} & \textbf{N} & \textbf{D} & \textbf{SD}\\
\hline
The text and visual elements of StatiCalcs are easy to read. & 146 & 59 & 0 & 0 & 0\\
The website works properly on different devices such as laptops and mobile phones. & 150 & 52 & 0 & 1 & 0\\
The layout makes important information easy to locate. & 154 & 48 & 0 & 0 & 0\\
The explanations and results provided by the website are easy to understand. & 144 & 51 & 0 & 0 & 0\\
Overall, StatiCalcs is accessible under different conditions. & 142 & 60 & 0 & 0 & 0\\
\hline
\end{tabular}
\end{table}

\begin{table}[H]
\centering
\caption{Weighted Mean of Students’ Perception on Accessibility}
\renewcommand{\arraystretch}{1.3}
\begin{tabular}{p{8cm}cc}
\hline
\textbf{Statement} & \textbf{Weighted Mean} & \textbf{Interpretation}\\
\hline
The text and visual elements of StatiCalcs are easy to read. & 4.66 & Strongly Agree\\
The website works properly on different devices such as laptops and mobile phones. & 4.66 & Strongly Agree\\
The layout makes important information easy to locate. & 4.69 & Strongly Agree\\
The explanations and results provided by the website are easy to understand. & 4.61 & Strongly Agree\\
Overall, StatiCalcs is accessible under different conditions. & 4.63 & Strongly Agree\\
\textbf{Overall Weighted Mean} & \textbf{4.65} & \textbf{Strongly Agree}\\
\hline
\end{tabular}
\end{table}

The results indicate that the student respondents strongly agreed that \textit{StatiCalcs} is accessible and functional across different devices and learning conditions. The statement, ``The layout makes important information easy to locate,'' obtained the highest weighted mean of 4.69, suggesting that the organization and presentation of information in the platform supported ease of use and navigation. On the other hand, the statement, ``The explanations and results provided by the website are easy to understand,'' obtained the lowest weighted mean of 4.61. Nevertheless, all indicators were interpreted as ``Strongly Agree,'' indicating that the respondents positively evaluated the accessibility features of the developed platform.

\section{Level of Perception of Students Regarding StatiCalcs in Terms of Satisfaction}

\begin{table}[H]
\centering
\caption{Frequency Distribution of Students’ Responses on Satisfaction}
\renewcommand{\arraystretch}{1.3}
\begin{tabular}{p{8cm}ccccc}
\hline
\textbf{Statement} & \textbf{SA} & \textbf{A} & \textbf{N} & \textbf{D} & \textbf{SD}\\
\hline
I am satisfied with my overall experience using StatiCalcs. & 157 & 44 & 0 & 0 & 0\\
StatiCalcs meets my expectations as a supplementary learning tool. & 155 & 44 & 0 & 0 & 0\\
The features of StatiCalcs are beneficial for practice and problem solving. & 170 & 33 & 0 & 0 & 0\\
I trust the accuracy of the calculations generated by StatiCalcs. & 140 & 50 & 0 & 0 & 0\\
Compared to textbooks, modules, and online videos, StatiCalcs provides a more convenient learning experience. & 133 & 46 & 0 & 2 & 0\\
\hline
\end{tabular}
\end{table}

\begin{table}[H]
\centering
\caption{Weighted Mean of Students’ Perception on Satisfaction}
\renewcommand{\arraystretch}{1.3}
\begin{tabular}{p{8cm}cc}
\hline
\textbf{Statement} & \textbf{Weighted Mean} & \textbf{Interpretation}\\
\hline
I am satisfied with my overall experience using StatiCalcs. & 4.69 & Strongly Agree\\
StatiCalcs meets my expectations as a supplementary learning tool. & 4.67 & Strongly Agree\\
The features of StatiCalcs are beneficial for practice and problem solving. & 4.76 & Strongly Agree\\
I trust the accuracy of the calculations generated by StatiCalcs. & 4.56 & Strongly Agree\\
Compared to textbooks, modules, and online videos, StatiCalcs provides a more convenient learning experience. & 4.46 & Strongly Agree\\
\textbf{Overall Weighted Mean} & \textbf{4.63} & \textbf{Strongly Agree}\\
\hline
\end{tabular}
\end{table}

The findings reveal that the student respondents were highly satisfied with the developed \textit{StatiCalcs} platform. The statement, ``The features of StatiCalcs are beneficial for practice and problem solving,'' obtained the highest weighted mean of 4.76, indicating that the respondents recognized the usefulness of the platform in solving engineering problems and reinforcing learning. Meanwhile, the statement comparing \textit{StatiCalcs} with traditional learning resources obtained the lowest weighted mean of 4.46. However, all indicators remained within the ``Strongly Agree'' interpretation range, suggesting that the respondents perceived \textit{StatiCalcs} as an effective and satisfactory supplementary learning tool for studying Statics of Rigid Bodies.

\section{Level of Preference of Students for Different Learning Resources in Studying Statics of Rigid Bodies}

\begin{table}[H]
\centering
\caption{Frequency Distribution of Students’ Preference for Learning Resources}
\renewcommand{\arraystretch}{1.3}
\begin{tabular}{p{6cm}ccccc}
\hline
\textbf{Learning Resource} & \textbf{MP} & \textbf{P} & \textbf{N} & \textbf{LP} & \textbf{LEP}\\
\hline
Textbooks & 128 & 68 & 14 & 2 & 0\\
Modules & 99 & 83 & 28 & 2 & 0\\
YouTube Videos & 107 & 73 & 28 & 3 & 1\\
Lecture Notes & 58 & 95 & 55 & 4 & 0\\
Online Engineering Tools & 88 & 50 & 62 & 10 & 2\\
StatiCalcs & 129 & 73 & 8 & 2 & 0\\
\hline
\end{tabular}
\end{table}

\begin{table}[H]
\centering
\caption{Weighted Mean of Students’ Preference for Learning Resources}
\renewcommand{\arraystretch}{1.3}
\begin{tabular}{p{6cm}cc}
\hline
\textbf{Learning Resource} & \textbf{Weighted Mean} & \textbf{Interpretation}\\
\hline
Textbooks & 4.52 & Most Preferred\\
Modules & 4.32 & Most Preferred\\
YouTube Videos & 4.33 & Most Preferred\\
Lecture Notes & 3.98 & Preferred\\
Online Engineering Tools & 4.00 & Preferred\\
StatiCalcs & 4.55 & Most Preferred\\
\hline
\end{tabular}
\end{table}

The findings reveal that \textit{StatiCalcs} was the most preferred learning resource among the student respondents, obtaining the highest weighted mean of 4.55, interpreted as ``Most Preferred.'' This was closely followed by textbooks, which obtained a weighted mean of 4.52. The results suggest that while students still value traditional learning resources, they also demonstrate strong acceptance of technology-supported learning tools for studying Statics of Rigid Bodies. Meanwhile, lecture notes obtained the lowest weighted mean of 3.98, interpreted as ``Preferred,'' indicating that students tend to favor more interactive and comprehensive learning resources.

\section{Level of Perception of Instructors Regarding StatiCalcs in Terms of Usability}

\begin{table}[H]
\centering
\caption{Frequency Distribution of Instructors’ Responses on Usability}
\renewcommand{\arraystretch}{1.3}
\begin{tabular}{p{8cm}ccccc}
\hline
\textbf{Statement} & \textbf{SA} & \textbf{A} & \textbf{N} & \textbf{D} & \textbf{SD}\\
\hline
The StatiCalcs website is easy to navigate. & 6 & 2 & 0 & 0 & 0\\
I can complete calculations quickly using the website. & 5 & 3 & 0 & 0 & 0\\
The website provides clear feedback after performing calculations. & 3 & 4 & 0 & 0 & 0\\
I can easily correct mistakes when entering inputs in the calculators. & 5 & 2 & 0 & 0 & 0\\
Overall, StatiCalcs is user-friendly. & 5 & 3 & 0 & 0 & 0\\
\hline
\end{tabular}
\end{table}

\begin{table}[H]
\centering
\caption{Weighted Mean of Instructors’ Perception on Usability}
\renewcommand{\arraystretch}{1.3}
\begin{tabular}{p{8cm}cc}
\hline
\textbf{Statement} & \textbf{Weighted Mean} & \textbf{Interpretation}\\
\hline
The StatiCalcs website is easy to navigate. & 4.75 & Strongly Agree\\
I can complete calculations quickly using the website. & 4.63 & Strongly Agree\\
The website provides clear feedback after performing calculations. & 4.25 & Strongly Agree\\
I can easily correct mistakes when entering inputs in the calculators. & 4.50 & Strongly Agree\\
Overall, StatiCalcs is user-friendly. & 4.63 & Strongly Agree\\
\textbf{Overall Weighted Mean} & \textbf{4.55} & \textbf{Strongly Agree}\\
\hline
\end{tabular}
\end{table}

The results indicate that the instructor respondents strongly agreed that \textit{StatiCalcs} is usable and easy to navigate. The statement, ``The StatiCalcs website is easy to navigate,'' obtained the highest weighted mean of 4.75, indicating a highly favorable evaluation of the platform’s interface and usability. Meanwhile, the statement, ``The website provides clear feedback after performing calculations,'' obtained the lowest weighted mean of 4.25. Nevertheless, all indicators were interpreted as ``Strongly Agree,'' demonstrating that the instructors positively perceived the usability of the developed learning platform.

\section{Level of Perception of Instructors Regarding StatiCalcs in Terms of Accessibility}

\begin{table}[H]
\centering
\caption{Frequency Distribution of Instructors’ Responses on Accessibility}
\renewcommand{\arraystretch}{1.3}
\begin{tabular}{p{8cm}ccccc}
\hline
\textbf{Statement} & \textbf{SA} & \textbf{A} & \textbf{N} & \textbf{D} & \textbf{SD}\\
\hline
The text and visual elements of StatiCalcs are easy to read. & 5 & 2 & 1 & 0 & 0\\
The website works properly on different devices such as laptops and mobile phones. & 5 & 3 & 0 & 0 & 0\\
The layout makes important information easy to locate. & 4 & 4 & 0 & 0 & 0\\
The explanations and results provided by the website are easy to understand. & 5 & 3 & 0 & 0 & 0\\
Overall, StatiCalcs is accessible under different conditions. & 2 & 5 & 1 & 0 & 0\\
\hline
\end{tabular}
\end{table}

\begin{table}[H]
\centering
\caption{Weighted Mean of Instructors’ Perception on Accessibility}
\renewcommand{\arraystretch}{1.3}
\begin{tabular}{p{8cm}cc}
\hline
\textbf{Statement} & \textbf{Weighted Mean} & \textbf{Interpretation}\\
\hline
The text and visual elements of StatiCalcs are easy to read. & 4.50 & Strongly Agree\\
The website works properly on different devices such as laptops and mobile phones. & 4.63 & Strongly Agree\\
The layout makes important information easy to locate. & 4.50 & Strongly Agree\\
The explanations and results provided by the website are easy to understand. & 4.63 & Strongly Agree\\
Overall, StatiCalcs is accessible under different conditions. & 4.13 & Agree\\
\textbf{Overall Weighted Mean} & \textbf{4.48} & \textbf{Strongly Agree}\\
\hline
\end{tabular}
\end{table}

The findings indicate that the instructor respondents generally perceived \textit{StatiCalcs} as accessible and functional across different conditions. The statements, ``The website works properly on different devices such as laptops and mobile phones'' and ``The explanations and results provided by the website are easy to understand,'' obtained the highest weighted means of 4.63. Meanwhile, the statement, ``Overall, StatiCalcs is accessible under different conditions,'' obtained the lowest weighted mean of 4.13, interpreted as ``Agree.'' Despite this, the overall weighted mean of 4.48 indicates that the instructors positively evaluated the accessibility of the developed web-based learning tool.

%===============================

\section{Level of Perception of Instructors Regarding StatiCalcs in Terms of Satisfaction}

\begin{table}[H]
\centering
\caption{Frequency Distribution of Instructors’ Responses on Satisfaction}
\renewcommand{\arraystretch}{1.3}
\begin{tabular}{p{8cm}ccccc}
\hline
\textbf{Statement} & \textbf{SA} & \textbf{A} & \textbf{N} & \textbf{D} & \textbf{SD}\\
\hline
I am satisfied with my overall experience using StatiCalcs. & 4 & 4 & 0 & 0 & 0\\
StatiCalcs meets my expectations as a supplementary learning tool. & 4 & 4 & 0 & 0 & 0\\
The features of StatiCalcs are beneficial for practice and problem solving. & 5 & 2 & 0 & 0 & 0\\
I trust the accuracy of the calculations generated by StatiCalcs. & 5 & 3 & 0 & 0 & 0\\
Compared to textbooks, modules, and online videos, StatiCalcs provides a more convenient learning experience. & 2 & 3 & 3 & 0 & 0\\
\hline
\end{tabular}
\end{table}

\begin{table}[H]
\centering
\caption{Weighted Mean of Instructors’ Perception on Satisfaction}
\renewcommand{\arraystretch}{1.3}
\begin{tabular}{p{8cm}cc}
\hline
\textbf{Statement} & \textbf{Weighted Mean} & \textbf{Interpretation}\\
\hline
I am satisfied with my overall experience using StatiCalcs. & 4.50 & Strongly Agree\\
StatiCalcs meets my expectations as a supplementary learning tool. & 4.50 & Strongly Agree\\
The features of StatiCalcs are beneficial for practice and problem solving. & 4.50 & Strongly Agree\\
I trust the accuracy of the calculations generated by StatiCalcs. & 4.63 & Strongly Agree\\
Compared to textbooks, modules, and online videos, StatiCalcs provides a more convenient learning experience. & 3.88 & Agree\\
\textbf{Overall Weighted Mean} & \textbf{4.40} & \textbf{Strongly Agree}\\
\hline
\end{tabular}
\end{table}

The findings reveal that the instructor respondents were generally satisfied with \textit{StatiCalcs} as a supplementary learning tool. The statement, ``I trust the accuracy of the calculations generated by StatiCalcs,'' obtained the highest weighted mean of 4.63, indicating confidence in the correctness and reliability of the platform’s outputs. Meanwhile, the statement comparing \textit{StatiCalcs} with textbooks, modules, and online videos obtained the lowest weighted mean of 3.88, interpreted as ``Agree.'' Nevertheless, the overall results indicate that instructors positively evaluated the developed platform in terms of satisfaction.

\section{Level of Preference of Instructors for Different Learning Resources in Studying Statics of Rigid Bodies}

\begin{table}[H]
\centering
\caption{Frequency Distribution of Instructors’ Preference for Learning Resources}
\renewcommand{\arraystretch}{1.3}
\begin{tabular}{p{6cm}ccccc}
\hline
\textbf{Learning Resource} & \textbf{MP} & \textbf{P} & \textbf{N} & \textbf{LP} & \textbf{LEP}\\
\hline
Textbooks & 6 & 1 & 0 & 1 & 0\\
Modules & 0 & 6 & 2 & 0 & 0\\
YouTube Videos & 0 & 3 & 3 & 1 & 1\\
Lecture Notes & 0 & 7 & 1 & 0 & 0\\
Online Engineering Tools & 1 & 4 & 3 & 0 & 0\\
StatiCalcs & 0 & 8 & 0 & 0 & 0\\
\hline
\end{tabular}
\end{table}

\begin{table}[H]
\centering
\caption{Weighted Mean of Instructors’ Preference for Learning Resources}
\renewcommand{\arraystretch}{1.3}
\begin{tabular}{p{6cm}cc}
\hline
\textbf{Learning Resource} & \textbf{Weighted Mean} & \textbf{Interpretation}\\
\hline
Textbooks & 4.50 & Most Preferred\\
Modules & 3.75 & Preferred\\
YouTube Videos & 3.00 & Neutral\\
Lecture Notes & 3.88 & Preferred\\
Online Engineering Tools & 3.75 & Preferred\\
StatiCalcs & 4.00 & Preferred\\
\hline
\end{tabular}
\end{table}

The results indicate that the instructor respondents still preferred traditional learning resources, particularly textbooks, which obtained the highest weighted mean of 4.50, interpreted as ``Most Preferred.'' \textit{StatiCalcs} obtained a weighted mean of 4.00, interpreted as ``Preferred,'' indicating positive acceptance of the developed web-based learning tool among instructors. Meanwhile, YouTube videos obtained the lowest weighted mean of 3.00, interpreted as ``Neutral.'' The findings suggest that instructors continue to value conventional academic resources while recognizing the usefulness of technology-supported learning tools.

\section{Significant Difference Between the Perceptions of Students and Instructors Regarding StatiCalcs}

\begin{table}[H]
\centering
\caption{Test of Significant Difference Between the Perceptions of Students and Instructors}
\small
\renewcommand{\arraystretch}{1.3}

\begin{tabular}{lccccc}
\hline
\textbf{Variable} &
\textbf{\makecell{Student\\Mean}} &
\textbf{\makecell{Instructor\\Mean}} &
\textbf{P-value} &
\textbf{Decision} &
\textbf{Interpretation}\\
\hline
Usability & 4.66 & 4.55 & 0.55 & Accept Ho & No Significant Difference\\
Accessibility & 4.65 & 4.48 & 0.62 & Accept Ho & No Significant Difference\\
Satisfaction & 4.63 & 4.40 & 0.27 & Accept Ho & No Significant Difference\\
\hline
\end{tabular}
\end{table}

The student respondents obtained weighted mean values of 4.66 for usability, 4.65 for accessibility, and 4.63 for satisfaction. These results indicate that the students positively perceived \textit{StatiCalcs} as a usable, accessible, and satisfactory web-based learning tool for studying Statics of Rigid Bodies. Among the evaluated variables, usability obtained the highest weighted mean, suggesting that the respondents found the platform easy to navigate and user-friendly.

On the other hand, the instructor respondents obtained weighted mean values of 4.55 for usability, 4.48 for accessibility, and 4.40 for satisfaction. The results likewise reflect positive perceptions toward the developed learning tool. Furthermore, all computed p-values were greater than the 0.05 level of significance; therefore, the null hypothesis was accepted. This indicates that there was no significant difference between the perceptions of students and instructors regarding the usability, accessibility, and satisfaction of \textit{StatiCalcs}.

\section{Respondents’ Comments and Suggestions Regarding StatiCalcs}

The respondents provided several comments, suggestions, and observations regarding the developed \textit{StatiCalcs} platform. Most responses expressed positive perceptions toward the usability, accessibility, and usefulness of the website as a supplementary learning tool for studying Statics of Rigid Bodies. Several respondents described the platform as user-friendly, easy to navigate, helpful, convenient, and effective for checking answers, reviewing lessons, and understanding engineering concepts. Many respondents also appreciated the integration of calculators, modules, formulas, and learning materials within a single platform. Some respondents expressed that the website would be highly beneficial for engineering students who struggle with Statics and related subjects.

Despite the positive feedback, respondents also provided several recommendations for future improvements of the platform. Common suggestions included the addition of step-by-step solutions, more examples and practice problems, formula displays, diagrams, and video tutorials to improve learning support and usability. Other respondents recommended improving navigation features by adding back and next buttons, enhancing formula visibility, implementing input validation, and strengthening mobile accessibility. Some respondents also suggested expanding the scope of the platform by incorporating additional engineering topics such as Dynamics of Rigid Bodies and Mechanics of Deformable Bodies, supporting multiple units of measurement, improving calculator functionality, and adding visualization features for better conceptual understanding. Overall, the comments and suggestions indicate that the respondents positively received the platform and recognized its strong potential for further enhancement and broader educational application.
